\chapter{Conclusioni}
\section{Limiti e prospettive future del lavoro}
Sebbene la versione attuale di TCFS funzioni nel quadro complessivo, non possiamo ignorare gli aspetti che sono rimasti incompleti o che funzionano solo parzialmente. L’algoritmo di cifratura AES GCM incontra notevoli difficoltà nel gestire blocchi di dimensioni superiori a 6 KB. Si è osservato che l’utilizzo della memoria aumenta drasticamente con l’avanzare del programma, nonostante la memoria residente di TCFS e quella virtuale rimangano invariate. È probabile che il problema sia la frammentazione della memoria dovuta a uno o più dei seguenti fattori: 
\begin{enumerate}
    \item \textbf{libssl}: Gestisce la memoria in modo particolare, cioè mantiene una lista interna di spazi di memoria liberi, invece di chiamare direttamente le funzioni malloc e free, il che può portare a una frammentazione interna della memoria. Questo perché la memoria libera può essere sparpagliata in diverse posizioni in porzioni piccole e non continue.
    \item \textbf{OS e compilatore}: Se il sistema operativo o il compilatore non liberano la memoria al momento dell’invocazione della funzione ‘free’ può verificarsi la frammentazione esterna. Questo accade quando la memoria libera è divisa in blocchi non contigui, impedendo l’allocazione di nuovi processi che richiedono blocchi di memoria più grandi.
    \item \textbf{Lazy memory allocation}: Questa tecnica, utilizzata per ritardare l’allocazione della memoria fino a quando non è effettivamente necessaria, può anch’essa contribuire alla frammentazione. Se la memoria non viene allocata immediatamente quando vengono richiesti numerosi byte, ciò può portare a rallentamenti e potenzialmente a frammentazione se le richieste di memoria successive richiedono blocchi di memoria contigui.
\end{enumerate}
Anche AES CTR può effettivamente riscontrare problemi quando si tratta di gestire blocchi di dimensioni considerevoli, in particolare quelli di dimensione pari o superiore a 32 MB. Questo è comprensibile, considerando che AES è un cifrario a blocchi con una dimensione di blocco fissa di 128 bit (16 byte). Anche se 32 MB possono sembrare una dimensione di blocco ragionevole, possono essere sufficientemente grandi da causare problemi di gestione della memoria, specialmente se la memoria non viene liberata in modo efficiente o se ci sono problemi di frammentazione della memoria. Tuttavia, è importante sottolineare che non esiste un motivo intrinseco per cui il plaintext debba essere un multiplo della dimensione del blocco del cifrario. Infatti, AES CTR esegue l’operazione XOR con il plaintext per ciascun blocco di output del cifrario. Ciò implica che il plaintext può avere qualsiasi lunghezza, e non necessariamente un multiplo della dimensione del blocco. Di conseguenza, AES CTR dovrebbe essere in grado di gestire messaggi di qualsiasi lunghezza. In conclusione, sebbene AES CTR possa riscontrare problemi con blocchi di 32 MB, ciò non dovrebbe limitare la sua utilità in molti contesti. Potrebbe essere necessario apportare ulteriori ottimizzazioni o considerare alternative se la dimensione del blocco diventa un problema.

Un altro problema significativo riguarda la gestione approssimativa degli IV. L’idea è di utilizzare un IV per la cifratura dei file e di utilizzare un IV diverso per ogni blocco del file. Tuttavia, l’IV è gestito correttamente solo per la cifratura della chiave del file. Questo ci porta al problema successivo: sarebbe necessario integrare una gestione della chiave dei file più sicura, simile a come è stata sviluppata la libreria key-utils. Questa gestione della chiave dei file non dovrebbe permettere di scrivere la chiave in un memory dump e dovrebbe sovrascrivere la chiave con 0 in memoria dopo l’utilizzo. Altre parti del progetto rimangono incomplete, il threshold sharing è stato implementato parzialmente, non è possibile recuperare lo share ottenuto da Shamir e non è stata implementata una utility per l'utente, in modo da agevolare l'utilizzo del FS. Inoltre sarebbe ideale creare un pacchetto per il FS e fornire in esso i binding per il comando mount in modo da poter utilizzare "mount.tcfs" per utilizzare il FS.

Ci si propone di risolvere e implementare tutti i cambiamenti necessari e implementare tutte le funzionalità mancanti in modo da rendere TCFS un FS completo e utilizzabile da tutti.